\documentclass[12pt]{article}
\usepackage[T1]{fontenc}
\usepackage[utf8]{inputenc}
\usepackage{lmodern}
\usepackage{textcomp}
\usepackage{enumitem}
\usepackage{geometry}
\geometry{margin=1in}
\begin{document}
\begin{center}
Answer Key - Health Sciences - K-12 Level Assessment 5 \
\small (Answers are ordered exactly as in the question paper.)
\end{center}

\section*{Multiple Choice}
\begin{enumerate}[label=\arabic*.]
\item \textbf{Answer:} D
\\ \textit{Options:} A) A reduction in lean body mass; B) A reduction in physical activity; C) A reduction in basal metabolic rate; D) All of the above
\\ \textit{Note:} All of the above factors, including decreased muscle mass, lower physical activity levels, and metabolic changes, collectively contribute to reduced energy expenditure in older adults.
\item \textbf{Answer:} D
\\ \textit{Options:} A) severe atopic dermatitis; B) enterocolitis; C) diarrhoea; D) all of the above
\\ \textit{Note:} The correct answer is "all of the above" because infants and children with a suspected food allergy, such as cow's milk allergy/intolerance, can exhibit a range of symptoms including gastrointestinal, skin, and respiratory reactions, indicating the diverse manifestations of food allergies.
\item \textbf{Answer:} D
\\ \textit{Options:} A) As simple Fe$_{2}$+ in the serum; B) Bound to albumin; C) Bound to ferritin; D) Bound to transferrin
\\ \textit{Note:} Iron is transported in the circulation bound to transferrin because this protein effectively binds and carries iron, ensuring its solubility and availability for metabolism at various sites in the body.
\item \textbf{Answer:} B
\\ \textit{Options:} A) Reduced lipolysis; B) Increased ketogenesis; C) Reduced gluconeogenesis; D) Reduced proteolysis
\\ \textit{Note:} Insulin deficiency leads to increased ketogenesis because without sufficient insulin, the body cannot effectively use glucose for energy, prompting it to break down fats for fuel, which produces ketones as a byproduct.
\item \textbf{Answer:} B
\\ \textit{Options:} A) Muscle contains less water than fat; B) Muscle contains more water than fat; C) Muscle weighs more than fat; D) Muscle weighs less than fat
\\ \textit{Note:} The correct answer is based on the fact that muscle tissue has a higher water content, which facilitates the conduction of electrical currents more efficiently than fat tissue, which contains less water.
\end{enumerate}

\section*{True / False}
\begin{enumerate}[label=\arabic*.]
\setcounter{enumi}{5}
\item \textbf{Answer:} True
\\ \textit{Options:} A) True; B) False
\\ \textit{Note:} The sarcomere is the basic functional unit of skeletal muscle tissue because it is the smallest contractile unit within the muscle fibers, responsible for the generation of muscle contractions through the interaction of actin and myosin filaments.
\item \textbf{Answer:} True
\\ \textit{Options:} A) True; B) False
\\ \textit{Note:} Leptin is a hormone produced by fat cells that signals the brain to reduce appetite and increase energy expenditure, while ghrelin, produced in the stomach, signals hunger, thus supporting the statement's accuracy.
\item \textbf{Answer:} False
\\ \textit{Options:} A) True; B) False
\\ \textit{Note:} The weighed diary is not the most commonly used dietary assessment method in epidemiological research; instead, methods like food frequency questionnaires or 24-hour dietary recalls are more prevalent due to their practicality and lower participant burden.
\item \textbf{Answer:} True
\\ \textit{Options:} A) True; B) False
\\ \textit{Note:} A food additive is considered safe when its Estimated Daily Intake (EDI) is less than its Acceptable Daily Intake (ADI) because the ADI represents the maximum amount that can be consumed daily over a lifetime without posing a significant risk to health.
\item \textbf{Answer:} True
\\ \textit{Options:} A) True; B) False
\\ \textit{Note:} The statement is true because during alcohol metabolism, ethanol is first converted into acetaldehyde by the enzyme alcohol dehydrogenase, and then acetaldehyde is further metabolized into acetate by acetaldehyde dehydrogenase, making acetaldehyde and acetate the primary metabolites of this process.
\end{enumerate}

\section*{Long Form Response}
\begin{enumerate}[label=\arabic*.]
\setcounter{enumi}{10}
\item \textbf{Answer:} Engel's law states that as people's income increases, the proportion of their income spent on food decreases. This phenomenon occurs because, with higher income levels, individuals and families can afford to allocate more money to other needs and desires beyond basic sustenance, such as education, healthcare, and leisure activities. For example, a family earning a low income might spend a large percentage of their budget on groceries to meet their basic needs, while a wealthier family might spend a smaller percentage on food, choosing to invest in other areas. This trend highlights how economic growth can influence consumer behavior and dietary choices, ultimately affecting food economics and health outcomes in society.
\\ \textit{Note:} The answer correctly illustrates Engel's law by demonstrating that as income rises, the percentage of income allocated to food decreases, allowing households to invest more in non-food necessities and discretionary spending.
\item \textbf{Answer:} Organic milk typically contains lower iodine levels compared to non-organic milk. This difference is primarily due to the feeding practices and farming methods used in organic dairy production, which may not include iodine-rich supplements that are more commonly utilized in non-organic farming. Iodine is essential for proper thyroid function and overall health, so lower iodine intake from organic milk could potentially lead to iodine deficiency in individuals who rely heavily on this source for their dietary needs. Therefore, while organic milk offers benefits such as fewer pesticides and antibiotics, consumers should be mindful of its reduced iodine content and consider incorporating other iodine-rich foods into their diet to maintain adequate levels for optimal health.
\\ \textit{Note:} The answer correctly identifies that organic milk usually has lower iodine levels due to the absence of iodine-rich supplements in organic farming practices, which is significant since iodine is crucial for thyroid function and overall health, potentially leading to deficiency in those who consume primarily organic milk.
\end{enumerate}

\vfill
\noindent\textit{End of Answer Key}
\end{document}